% Generated draft: requires mathematical and source review.
\documentclass[11pt]{article}
\usepackage[margin=25mm]{geometry}
\usepackage{fontspec}
\setmainfont{Latin Modern Roman}
\usepackage{amsmath,amssymb,mathtools,unicode-math}
\setmathfont{Latin Modern Math}
\usepackage{enumitem,needspace,xurl,hyperref}
\setlength{\parindent}{0pt}
\setlength{\parskip}{5pt}
\setlength{\emergencystretch}{3em}
\newcommand{\N}{\mathbb{N}}
\newcommand{\Z}{\mathbb{Z}}
\newcommand{\Q}{\mathbb{Q}}
\newcommand{\R}{\mathbb{R}}
\newcommand{\C}{\mathbb{C}}
\newcommand{\F}{\mathbb{F}}
\newcommand{\A}{\mathbb{A}}
\newcommand{\PP}{\mathbb{P}}
\newcommand{\E}{\mathbb{E}}
\newcommand{\eps}{\varepsilon}
\newcommand{\dd}{\,\mathrm{d}}
\begin{document}
\section*{UnsolvedMath ID: 30000051}
\textbf{Source problem code:} catalogue entry
\subsection{Definitions and mathematical statement}
For integers $n,d\ge1$ and an $n$-point set $M=\{x_1,\ldots,x_n\}\subset[0,1]^d$, define its star discrepancy by
\[
 D_n^*(M)=\sup_{t\in[0,1]^d}
 \left|\frac1n\#\left(M\cap\prod_{j=1}^d[0,t_j)\right)
                     -\prod_{j=1}^dt_j\right|.
\]
The requested construction has a single absolute constant $C$ independent of both $n$ and $d$ such that
\[
 D_n^*(M)\le C\sqrt{d/n}.
\]
The main question asks for an explicit construction, in time polynomial in $n$ and $d$, of such a point set for every pair $(n,d)$. The weaker alternative asks for an explicitly described small finite candidate family
\[
 S_{n,d}\subset[0,1]^{dn}
\]
guaranteed to contain an ordered configuration of $n$ distinct points whose associated set has the same bound. The cube of dimension $dn$ encodes all $nd$ coordinates; it is not the point domain itself.

The source does not quantify ``small'', prescribe a seed length or success fraction in the family, or choose a bit-complexity versus real-arithmetic model for coordinate output. These must be specified for a sharp algorithmic theorem. In particular, an unspecified singleton containing an existentially chosen good configuration would not satisfy the intended explicit or reduced-randomness goal. A guarantee that a family contains one good configuration is also distinct from an efficient procedure selecting that member.

The literature assessment checked on 21 August 2026 says the existential $O(\sqrt{d/n})$ bound is classical and that explicit and pseudorandom constructions give substantial progress, while a small finite sample space or fully explicit polynomial-time construction at the stated high-dimensional goal remains open in that assessment. The question concerns uniform dependence on dimension, so a fixed-dimension discrepancy bound with an uncontrolled dimension-dependent constant is insufficient. These status claims are dated curation requiring independent verification.
\subsection{Short English statement}
Construct high-dimensional point sets with star discrepancy at most an absolute constant times $\sqrt{d/n}$ in polynomial time, or explicitly give a small candidate family containing one. The second option reduces randomness but does not by itself select a good candidate.
\subsection{Sources}
\begin{itemize}
\item \emph{Discrepancy Theory and Its Applications} (2004), Oberwolfach Reports, DOI 10.4171/owr/2004/13, pp. 673--722, \url{https://doi.org/10.4171/owr/2004/13}.
\item Literature assessment checked 21 August 2026: ``Pseudorandom generators and discrepancy,'' \url{https://arxiv.org/abs/2004.14158}; ``Explicit low-discrepancy constructions,'' \url{https://arxiv.org/abs/1308.4252}; ``OWR 13/2004,'' \url{https://doi.org/10.4171/OWR/2004/13}.
\item Local UnsolvedMath registry, permanent ID $30000051$, exact source code \texttt{catalogue entry}, created 25 August 2026 and updated 16 September 2026. The clean statement is labelled corrected and verified against the official report, restoring the quantitative target and polynomial-time alternative.
\end{itemize}
\end{document}
